% GENERATED FILE. Edit canonical lessons, then run npm run generate:books.
% canonical-source-hash: fnv1a64:9bd7bdd36dae9b15
% canonical-lessons: ZH-C33-nan, ZH-W33-nan, ZH-C33-nanren, ZH-C33-nuren, ZH-C33-meimei, ZH-C33-haizi

\chapter{The Rest of the Family}
\label{ch:zh-the-rest-of-the-family}
\hlchaptermodality{33}

\begin{chapteropening}
\textbf{By the end of this chapter:} \emph{I can talk about men, women, sisters and children.}

Complete the last lesson of chapter 33: I can talk about men, women, sisters and children.
\end{chapteropening}

\section[\texorpdfstring{\zh{男} (nán)}{男 (nán)}]{\zh{男} (nán) — male}

\label{lesson:ZH-C33-nan}



\begin{quote}
Before the new one: say the Mandarin for elder sister, then the Mandarin for younger brother.
\end{quote}

\subsection*{You'll want to know: \zh{男}}
\textbf{\zh{男}} — \emph{nán} — \textquotedblleft{}male\textquotedblright{}.

\begin{quote}
\textbf{\zh{我有一个男同学}} \emph{wǒ yǒu yí ge nán tóngxué} — One of my classmates is a boy.
\end{quote}

The new character, in pieces: \zh{男} — a field on top, built like \zh{日} with a cross inside, and below it a two-stroke piece meaning strength. The next lesson puts it in your hand.

\begin{cousinweb}
\textbf{\zh{男}} (\emph{nán}) is inherited from Middle Chinese \emph{nom} (Guangyun, in Baxter's transcription), glossed 'a man'.
\end{cousinweb}

\begin{grammarlens}[title={What you've built}]
One more word about people.
\end{grammarlens}

\subsection*{Your turn}
\begin{itemize}
\raggedright
  \item \emph{Say it:} \emph{nán}
  \item \emph{Say it:} \emph{nán}, once more
  \item \emph{Say it:} \emph{wǒ yǒu yí ge nán tóngxué} — One of my classmates is a boy.
  \item \emph{Recall:} say the Mandarin for elder sister, then the Mandarin for younger brother, then say \emph{nán} again
\end{itemize}

\begin{tcolorbox}[breakable,colback=teal!4,colframe=teal!35!black,title={Before you move on}]
What does \zh{男} mean? (\textquotedblleft{}Male\textquotedblright{}.) Say it once more.
\end{tcolorbox}
\section[\texorpdfstring{\zh{男} (nán)}{男 (nán)}]{\zh{男} — a field above, strength below}

\label{lesson:ZH-W33-nan}



\begin{quote}
The word first, then the character inside it.

\begin{itemize}
\raggedright
  \item \emph{Recall:} say \emph{nán} — \textbf{R1}, one lesson back, before it decays
\end{itemize}

\begin{quote}
\zh{日} \zh{口}
\end{quote}

\emph{Read it:} each one aloud, with nothing said first — \emph{rì}, \emph{kǒu}
\end{quote}

\subsection*{Script — a field above, strength below}
\hlblockfigure{\detokenize{figures/ZH-W33-nan-filmstrip.pdf}}{How \zh{男} is written, stroke by stroke}

\begin{quote}
\zh{男}
\end{quote}

\textbf{Two pieces, top before bottom.}

\textbf{The top piece, a field, five strokes.} A vertical down the left. A horizontal across the top that turns down the right side. The short horizontal across the middle. The short vertical down the middle. Then the horizontal that closes the bottom. It is \zh{日} with a cross inside instead of a single bar.

\textbf{The bottom piece, strength, two strokes.} A horizontal that turns down and hooks to the left, then a long left-falling stroke through it.

5 + 2 = seven.

\subsection*{Your turn}
\begin{itemize}
\raggedright
  \item \emph{Write it:} \zh{男}, all seven, with the model in front of you
  \item \emph{Cover:} write it again from memory, and count the pieces, not the strokes
\end{itemize}

\begin{tcolorbox}[breakable,colback=teal!4,colframe=teal!35!black,title={Before you move on}]
How many pieces in \zh{男}? (\textbf{Two} — the field, then strength.) Which character you know does the field look like? (\textbf{\zh{日}}, with a cross inside.) Which stroke closes the field? (\textbf{The bottom horizontal, the fifth}.)

Source: \href{https://raw.githubusercontent.com/chanind/hanzi-writer-data/68d10a4b21150cae5e1ebbd223eed289cf32d90c/data/\%E7\%94\%B7.json}{Hanzi Writer Data \zh{男}, pinned stroke paths and medians}.
\end{tcolorbox}
\section[\texorpdfstring{\zh{男人} (nánrén)}{男人 (nánrén)}]{\zh{男人} (nánrén) — man}

\label{lesson:ZH-C33-nanren}



\begin{quote}
Before the new one: say the Mandarin for younger brother, then the Mandarin for male.
\end{quote}

\subsection*{You'll want to know: \zh{男人}}
\textbf{\zh{男人}} — \emph{nánrén} — \textquotedblleft{}man\textquotedblright{}.

\begin{quote}
\textbf{\zh{我家有三个男人}} \emph{wǒ jiā yǒu sān ge nánrén} — There are three men in my family.
\end{quote}

\begin{cousinweb}
\textbf{\zh{男人}} (\emph{nánrén}) is built from \zh{男} 'male' and \zh{人} 'person', both inherited from Middle Chinese, \emph{nom} and \emph{nyin} (Guangyun, in Baxter's transcription).
\end{cousinweb}

\begin{grammarlens}[title={What you've built}]
One more word about people.
\end{grammarlens}

\subsection*{Your turn}
\begin{itemize}
\raggedright
  \item \emph{Say it:} \emph{nánrén}
  \item \emph{Say it:} \emph{nánrén}, once more
  \item \emph{Say it:} \emph{wǒ jiā yǒu sān ge nánrén} — There are three men in my family.
  \item \emph{Recall:} say the Mandarin for younger brother, then the Mandarin for male, then say \emph{nánrén} again
\end{itemize}

\begin{tcolorbox}[breakable,colback=teal!4,colframe=teal!35!black,title={Before you move on}]
What does \zh{男人} mean? (\textquotedblleft{}Man\textquotedblright{}.) Say it once more.
\end{tcolorbox}
\section[\texorpdfstring{\zh{女人} (nǚrén)}{女人 (nǚrén)}]{\zh{女人} (nǚrén) — woman}

\label{lesson:ZH-C33-nuren}



\begin{quote}
Before the new one: say the Mandarin for male, then the Mandarin for man.
\end{quote}

\subsection*{You'll want to know: \zh{女人}}
\textbf{\zh{女人}} — \emph{nǚrén} — \textquotedblleft{}woman\textquotedblright{}.

\begin{quote}
\textbf{\zh{我家有四个女人}} \emph{wǒ jiā yǒu sì ge nǚrén} — There are four women in my family.
\end{quote}

\begin{cousinweb}
\textbf{\zh{女人}} (\emph{nǚrén}) is built from \zh{女} 'female' and \zh{人} 'person', both inherited from Middle Chinese, \emph{nrjoX} and \emph{nyin} (Guangyun, in Baxter's transcription).
\end{cousinweb}

\begin{grammarlens}[title={What you've built}]
One more word about people.
\end{grammarlens}

\subsection*{Your turn}
\begin{itemize}
\raggedright
  \item \emph{Say it:} \emph{nǚrén}
  \item \emph{Say it:} \emph{nǚrén}, once more
  \item \emph{Say it:} \emph{wǒ jiā yǒu sì ge nǚrén} — There are four women in my family.
  \item \emph{Recall:} say the Mandarin for male, then the Mandarin for man, then say \emph{nǚrén} again
\end{itemize}

\begin{tcolorbox}[breakable,colback=teal!4,colframe=teal!35!black,title={Before you move on}]
What does \zh{女人} mean? (\textquotedblleft{}Woman\textquotedblright{}.) Say it once more.
\end{tcolorbox}
\section[mèimei]{mèimei — younger sister}

\label{lesson:ZH-C33-meimei}



\begin{quote}
Before the new one: say the Mandarin for man, then the Mandarin for woman.
\end{quote}

\subsection*{You'll want to know: mèimei}
\emph{mèimei} — \textquotedblleft{}younger sister\textquotedblright{}.

\begin{quote}
\emph{wǒ mèimei shì xiǎoxuésheng} — My little sister is at primary school.
\end{quote}

Its characters come later; for now it is a word for the ear.

\begin{cousinweb}
\emph{mèimei} is inherited from Middle Chinese \emph{mwojH} (Guangyun, in Baxter's transcription), where it is glossed with the word for sisters.
\end{cousinweb}

\begin{grammarlens}[title={What you've built}]
One more word about people.
\end{grammarlens}

\subsection*{Your turn}
\begin{itemize}
\raggedright
  \item \emph{Say it:} \emph{mèimei}
  \item \emph{Say it:} \emph{mèimei}, once more
  \item \emph{Say it:} \emph{wǒ mèimei shì xiǎoxuésheng} — My little sister is at primary school.
  \item \emph{Recall:} say the Mandarin for man, then the Mandarin for woman, then say \emph{mèimei} again
\end{itemize}

\begin{tcolorbox}[breakable,colback=teal!4,colframe=teal!35!black,title={Before you move on}]
What does mèimei mean? (\textquotedblleft{}Younger sister\textquotedblright{}.) Say it once more.
\end{tcolorbox}
\section[háizi]{háizi — child}

\label{lesson:ZH-C33-haizi}



\begin{quote}
Before the new one: say the Mandarin for woman, then the Mandarin for younger sister.
\end{quote}

\subsection*{You'll want to know: háizi}
\emph{háizi} — \textquotedblleft{}child\textquotedblright{}.

\begin{quote}
\emph{tāmen yǒu sān ge háizi} — They have three children.
\end{quote}

Its characters come later; for now it is a word for the ear.

\begin{cousinweb}
\emph{háizi} is built from hái 'child' and the noun ending zi. The first is inherited from Middle Chinese \emph{hoj} (Guangyun, in Baxter's transcription), glossed 'a newborn child'; the ending is \emph{tsiX}, 'child', the same piece as the one in \zh{儿子}.
\end{cousinweb}

\begin{grammarlens}[title={What you've built}]
One more word about people.
\end{grammarlens}

\subsection*{Your turn}
\begin{itemize}
\raggedright
  \item \emph{Say it:} \emph{háizi}
  \item \emph{Say it:} \emph{háizi}, once more
  \item \emph{Say it:} \emph{tāmen yǒu sān ge háizi} — They have three children.
  \item \emph{Recall:} say the Mandarin for woman, then the Mandarin for younger sister, then say \emph{háizi} again
\end{itemize}

\begin{tcolorbox}[breakable,colback=teal!4,colframe=teal!35!black,title={Before you move on}]
What does háizi mean? (\textquotedblleft{}Child\textquotedblright{}.) Say it once more.
\end{tcolorbox}
